% 第4章 观测增强与 HIL 学习的鲁棒任务策略

\chapter{观测增强与 HIL 学习的鲁棒任务策略}\label{ch:task_policy}

\section{方法设计动机}\label{sec:task_method_motivation}

\cref{ch:problem_formulation} 从物理因果链与 POMDP 形式化两个层面刻画了编码器偏差问题的结构，并通过可观测性分析指出"信息上能辨识 $\vect{b}$"的充分条件。然而理论分析并不直接告诉我们"应当设计什么样的策略学习方法"。本节的目的是承接\cref{ch:problem_formulation} 的结论，将它们转化为三个具体的方法设计选择，并简要说明每个选择的依据与对应章节。

\paragraph{设计选择一：观测空间增强} 根据\cref{prop:non_identifiable}，仅依赖基础本体感觉观测——其中关节角读数 $\vect{q}+\vect{b}$ 以及由其经正运动学得到的 TCP 位置 $\text{FK}(\vect{q}+\vect{b})$ 均被偏差污染，而关节速度（编码器读数对时间求导，静态偏差 $\dot{\vect{b}}=\vect{0}$ 被消去）与夹爪开度（来自控制指令而非编码器）不受偏差影响——\emph{无法}区分不同的 $(\vect{q}_{\text{true}}, \vect{b})$ 对；而\cref{prop:identifiable} 表明，引入一个独立于编码器的外部位置观测 $\vect{o}_{\text{real}}$ 可以打破 $(\vect{q}_{\text{true}}, \vect{b})$ 之间的平移对称性，使系统近似可辨识。因此本文第一个设计选择是\textbf{在观测中显式加入外部位置信号}，在原有基础观测之上扩展为 24 维观测空间。具体实现见\cref{sec:obs_design}。

\paragraph{设计选择二：随机偏差训练} \cref{subsec:identifiability_vs_solvability} 进一步指出，即便信息充分，最优策略仍然是\emph{以偏差为条件}的函数：对不同的 $\vect{b}$，$\pi^\star$ 必须输出不同的补偿动作（见\cref{eq:conditional_policy}）。若训练过程仅暴露单一固定 $\vect{b}$，策略只能在该点附近学到"单点条件化"——在偏差分布的其他位置（尤其是 $\vect{b}=\vect{0}$）会产生过补偿现象。因此本文第二个设计选择是\textbf{按 episode 从偏差分布 $P(\vect{b})$ 中随机采样进行训练}，让策略暴露于整个偏差分布。这一思路与 Tobin 等人在视觉 sim-to-real 中采用的域随机化~\cite{tobin2017domain} 在原则上同源——两者都是"训练时让策略见过未来部署环境的分布"。\cref{sec:random_bias} 将具体描述随机偏差训练，\cref{sec:exp_h2} 的固定偏差基线实验将给出单点条件化过补偿的直接实证。

\paragraph{设计选择三：HIL + RLPD 真机训练} 前两个选择回答了"观测什么"和"训练什么分布"，但仍然留下执行层面的问题：\emph{在哪里训练}？本文的任务策略必须在\textbf{真机上}训练而非纯仿真 + 域随机化直接迁移，这出于两个现实约束。(i)~\textbf{视觉 sim-to-real gap}：任务策略的观测除了 24 维状态还包含双目 RGB 图像，由 ResNet 编码器处理~\cite{he2016deep}；仿真渲染图像与真实相机图像在纹理、光照与噪声分布上的差距难以仅通过 DR 覆盖（这一观察与\cref{ch:backup_policy} 备份策略可以纯仿真训练的情形形成对照，详见\cref{ch:backup_policy} 5.1 节与\cref{ch:related_work}）。(ii)~\textbf{稀疏奖励下的探索困难}：本文采用成功即 1、否则为 0 的稀疏奖励（见\cref{sec:task_reward}），策略难以从零开始探索出成功轨迹。为此本文采用 Ball 等人提出的 RLPD 框架~\cite{ball2023rlpd}：将 50 条人类演示 episode 放入离线缓冲区、与在线采集的数据混合采样，在 SAC 基础上提供稠密的引导信号。训练过程中进一步引入\emph{人在回路干预}——这一思路可追溯到 Ross 等人提出的 DAgger~\cite{ross2011dagger}：训练者通过遥操作设备介入失败轨迹，成功的干预被加入离线缓冲区形成与演示等同地位的监督信号。实验表明，在稀疏奖励与随机偏差的双重设定下，HIL 干预在\emph{仿真和真机上均是必要组成}——即便在仿真环境中，仅依靠 demo bootstrap + 纯 online SAC 也难以稳定收敛，需要人类干预持续提供纠错梯度（\cref{sec:exp_h1} 的残差 RL 对照实验将从另一个角度佐证这一点）。训练框架、网络架构与超参数分别见\cref{sec:rlpd,sec:network}，真机部署细节参见\cref{ch:real_robot}。

\paragraph{小结} 上述三个设计选择可以简明地映射为\cref{tab:method_mapping} 所示的"理论-约束-方法"对应关系。后续各节将按此顺序展开具体实现。

\begin{table}[htbp]
  \centering
  \caption{从\cref{ch:problem_formulation} 的理论分析与现实约束到本章方法设计的映射}
  \label{tab:method_mapping}
  \footnotesize
  \begin{tabular}{@{}llll@{}}
    \toprule
    来源 & 具体依据 & 方法设计 & 展开章节 \\
    \midrule
    理论结论 & 命题~\ref{prop:identifiable}（可辨识性） & 观测空间增强（24D） & \cref{sec:obs_design} \\
    理论结论 & \cref{subsec:identifiability_vs_solvability}（条件补偿） & 随机偏差训练 & \cref{sec:random_bias} \\
    现实约束 & 视觉 sim-to-real gap & 真机训练 & \cref{sec:rlpd} \\
    现实约束 & 稀疏奖励探索困难 & RLPD + HIL 干预 & \cref{sec:rlpd} \\
    \bottomrule
  \end{tabular}
\end{table}

\section{观测空间设计}\label{sec:obs_design}

\subsection{观测向量构成}\label{subsec:obs_vector}

任务策略的低维状态观测向量 $\vect{o}_t \in \mathbb{R}^{24}$ 由六组分量顺序拼接而成（\cref{tab:obs_24d}）：前 18 维来自机器人本体感觉通路（编码器读数与控制指令），后 6 维来自独立于编码器的外部观测通路。$\vect{o}_\text{real}$ 的噪声建模与 TCP 定位方案的抽象化处理见\cref{subsec:obs_image_tcp}。

\cref{tab:obs_24d} 同时标注了各分量的两个正交性质——\emph{是否受偏差 $\vect{b}$ 污染}与\emph{对任务目标的相关性}。后者分为三类："直接相关"（$\vect{o}_\text{block}$、$\vect{o}_\text{real}$ 与 $\text{FK}(\vect{q}+\vect{b})$）直接参与"靠近并抓取方块"的决策，"间接相关"的关节角 $\vect{q}+\vect{b}$ 仅在关节极限或几何约束触发时被利用，"结构性"的关节速度与夹爪开度主要服务于控制闭环的平滑性。这一分类揭示出本节的核心设计思路：三个直接任务相关的位置量中，$\text{FK}(\vect{q}+\vect{b})$ 与 $\vect{q}+\vect{b}$ 受污染、$\vect{o}_\text{real}$ 未受污染——正是这组\emph{直接相关且彼此不一致}的位置信号为策略隐式辨识 $\vect{b}$ 提供了信息来源。

\begin{table}[htbp]
  \centering
  \caption{任务策略的 24 维状态观测构成}
  \label{tab:obs_24d}
  \small
  \begin{tabular}{@{}clclcc@{}}
    \toprule
    索引 & 分量 & 维度 & 来源 & 受偏差污染 & 任务相关 \\
    \midrule
    $[0, 6]$   & 关节角 $\vect{q}+\vect{b}$ & 7 & 编码器读数 & 是 & 间接 \\
    $[7, 13]$  & 关节速度 $\dot{\vect{q}}$ & 7 & 编码器读数对时间差分 & 否 & 结构性 \\
    $14$       & 夹爪开度 & 1 & 控制指令 & 否 & 结构性 \\
    $[15, 17]$ & TCP 位置 $\text{FK}(\vect{q}+\vect{b})$ & 3 & 有偏正运动学 & 是 & 直接 \\
    $[18, 20]$ & 方块位置 $\vect{o}_\text{block}$ & 3 & 场景外部相机 & 否 & 直接 \\
    $[21, 23]$ & 真实 TCP $\vect{o}_\text{real}$ & 3 & 外部定位 + 5\,mm 噪声 & 否 & 直接 \\
    \bottomrule
  \end{tabular}
\end{table}

\subsection{信息通路}\label{subsec:obs_dataflow}

\cref{fig:obs_dataflow} 示意了 24 维观测在两条独立通路上的流动（污染标记见图中红色标签）。根据\cref{prop:identifiable}，两路联合观测已满足近似可辨识性的充分条件——策略原则上可通过比较 $\vect{q}+\vect{b}$ 与 $\vect{o}_\text{real}$ 的不一致隐式恢复 $\vect{b}$，并据此输出\cref{subsec:identifiability_vs_solvability} 所要求的条件补偿动作。本文不对估计结构施加显式约束（如卡尔曼滤波或专门的 $\vect{b}$ 预测头），而交由端到端优化自行组织。

\begin{figure}[htbp]
  \centering
  \begin{tikzpicture}[
    >={Stealth[length=2mm]},
    font=\small,
    grp/.style={draw, rounded corners=4pt, inner sep=8pt, align=left},
    policy/.style={draw, fill=gray!15, rounded corners=3pt, inner sep=6pt, align=center, minimum width=4cm},
  ]
    \node[grp, label={[font=\small\bfseries]above:编码器路径 (18D)}] (enc) {%
      $\vect{q}+\vect{b}$~关节角 (7D)~\textcolor{red!70!black}{\scriptsize [污染]}\\
      $\dot{\vect{q}}$~关节速度 (7D)\\
      夹爪开度 (1D)\\
      $\text{FK}(\vect{q}+\vect{b})$~TCP 位置 (3D)~\textcolor{red!70!black}{\scriptsize [污染]}%
    };
    \node[grp, right=1.6cm of enc, label={[font=\small\bfseries]above:外部观测路径 (6D)}] (ext) {%
      $\vect{o}_\text{block}$~方块位置 (3D)\\
      $\vect{o}_\text{real}$~真实 TCP (3D)%
    };

    \node[policy, below=1.2cm of $(enc.south)!0.5!(ext.south)$] (pi) {策略 $\pi_\theta$};

    \draw[->] (enc.south) -- (enc.south |- pi.north);
    \draw[->] (ext.south) -- (ext.south |- pi.north);
    \draw[->] (pi.south) -- ++(0,-0.55) node[right, font=\small] {动作 $\vect{a}_t$};
  \end{tikzpicture}
  \caption{24 维观测空间的信息通路示意。编码器路径（左）中的关节角与由其导出的 TCP 位置被偏差 $\vect{b}$ 污染，关节速度与夹爪开度不受影响；外部观测路径（右）与编码器完全解耦。两路信号在策略网络处汇合，策略原则上可通过两者之间的不一致隐式恢复 $\vect{b}$。}
  \label{fig:obs_dataflow}
\end{figure}

\subsection{为何保留被污染的本体感觉}\label{subsec:obs_keep_contam}

一个自然的疑问是：既然 $\vect{q}+\vect{b}$ 与 $\text{FK}(\vect{q}+\vect{b})$ 已被偏差污染，为何不直接剔除、仅以 $\vect{o}_\text{real}$ 作为 TCP 的唯一来源？本文选择保留两者，基于两点考虑。

\textbf{信息完整性。}按\cref{subsec:identifiability_vs_solvability} 的分析，策略能够辨识 $\vect{b}$ 的前提是同时持有"被污染的编码器信号"与"独立的外部位置信号"——两者之\emph{差}才承载偏差信息。若剔除前者，策略仅剩 $\vect{o}_\text{real}$ 一路，反而失去推断 $\vect{b}$ 所需的另一半观测。

\textbf{运行时容错。}真机外部定位系统（ArUco/AprilTag、OptiTrack 等）可能出现短暂遮挡、标记丢失或通讯抖动；保留本体感觉通路意味着当 $\vect{o}_\text{real}$ 暂时不可用时，策略仍能退化到"仅依靠本体感觉"的带偏差行为，而非彻底失效。

\subsection{图像观测与 TCP 定位的通用化处理}\label{subsec:obs_image_tcp}

\paragraph{图像观测。}任务策略同时接收前视相机（场景全局视角）与腕部相机（夹爪局部视角）两路 RGB 图像，经 ResNet 编码~\cite{he2016deep} 后与 24 维状态观测拼接作为策略输入（视觉分支结构见\cref{sec:network}）。图像在仿真渲染与真实相机之间存在显著的纹理、光照与噪声差距，这正是\cref{sec:task_method_motivation} 论证"任务策略必须在真机上训练"的直接依据。

\paragraph{\texttt{real\_tcp} 的噪声建模。}仿真训练中，$\vect{o}_\text{real}$ 由 MuJoCo 场景真值叠加 $\mathcal{N}(\vect{0}, (0.005\,\text{m})^2\mathbf{I})$ 高斯噪声得到，以模拟真机外部定位系统的典型单帧误差量级，降低 sim-to-real 迁移时因未建模噪声引起的性能衰减。

\paragraph{TCP 定位方案的抽象化。}本章不将 $\vect{o}_\text{real}$ 绑定到具体实现——可选方案包括 ArUco/AprilTag 视觉标记、OptiTrack 光学动捕等，只要该通路满足三条最小要求：独立于关节编码器、能以米制给出 TCP 位置、典型噪声量级与 5\,mm 可比。具体选型与标定细节见\cref{ch:real_robot}。

\section{随机偏差训练}\label{sec:random_bias}

% 动机：将 $\vect{b}$ 视为环境参数 → 域随机化的一种形式
% 训练时 $\vect{b} \sim U[0, b_{\max}]$ per episode，episode内恒定
%
% 与固定偏差训练的对比：
% - 固定 $b=0.2$：过拟合到特定 $\vect{b}$ → 零偏差时过补偿
% - 随机 $b \in [0, 0.25]$：学到 $\vect{b}$-条件策略 $\pi(\vect{a}|\vect{o}, \Delta\vect{x})$
%
% 与鲁棒RL的关系：
% - 不是最坏情况（鲁棒MDP），而是偏差分布上的平均性能
% - 随机化迫使策略利用 $\Delta\vect{x}$ 信号推断 $\vect{b}$

\section{RLPD训练框架}\label{sec:rlpd}

% SAC（Soft Actor-Critic）：适合连续动作空间
%
% RLPD = SAC + 离线演示缓冲区 + 在线缓冲区
% - 离线缓冲区：50个人类演示episode
% - 在线缓冲区：训练中采集（容量100K）
% - 预热：500步纯离线训练，之后混合采样
% - 回放比例：50%离线 + 50%在线
%
% 人在回路：
% - 键盘干预（在线采集过程中）
% - 成功的干预加入离线缓冲区
% - 失败的干预丢弃
%
% 分布式架构：Actor ←gRPC→ Learner

\section{网络架构}\label{sec:network}

% 状态编码器：Linear(24, 256) + LayerNorm + Tanh
% 图像编码器：ResNet10（冻结）处理前视/腕部相机
% Actor：[256, 256] MLP → 6D连续 + 1D离散夹爪
% Critic：2× [256, 256] MLP（双Q学习）
% 温度：可学习的 $\log\alpha$，自动调节
%
% SAC损失函数：
% $L_Q = \expect[(Q(\vect{s},\vect{a}) - (r + \gamma \min_i Q_{\text{target},i}(\vect{s}', \vect{a}')))^2]$
% $L_\pi = \expect[\alpha \log \pi(\vect{a}|\vect{s}) - Q(\vect{s}, \vect{a})]$
% $L_\alpha = -\expect[\alpha (\log \pi(\vect{a}|\vect{s}) + \mathcal{H}_{\text{target}})]$
%
% 关键超参数：batch_size=256, utd_ratio=2, discount=0.97, lr=3e-4

\section{奖励设计}\label{sec:task_reward}

% 抓取任务稀疏奖励：
% $r = 1.0$ if success else $0.0$
% success = (距离 < 5cm) AND (抬升高度 > 20cm)
%
% 稀疏奖励可行的原因：
% - RLPD提供了演示引导的探索
% - 人类干预引导早期训练
% - 密集奖励在偏差下容易产生reward hacking
